\chapter{分子之间也会彼此吸引}

\begin{kaichang}
做饭的时候撕一段保鲜膜，往碗口上一蒙，用手抹一圈——它居然自己贴住了，不用胶水，不用胶带。更奇怪的是，你随时能把它撕下来，再贴回去，再撕下来。它到底靠什么粘住？又为什么粘得这么“客气”，一撕就开？

先别急着回答，再看看这几件事：清晨草叶上的露珠，为什么总是一颗一颗的小圆球，而不是摊成薄薄一片？壁虎能倒挂在天花板上散步，它的脚底没有吸盘，也没有胶水，靠什么对抗重力？在房间里喷一下香水，几分钟后整个屋子都闻得到，而瓶子里的液体一天比一天少——跑掉的香水去了哪里？

四件看起来毫不相干的小事。你先猜：它们背后是同一种原因，还是四个互不相干的原因？把猜测记下来，这一章结束时回来对答案。
\end{kaichang}

\section{粘住、聚成球、悄悄逃走}

先把眼睛能看到的现象摆出来。

第一类是“粘”：保鲜膜贴上碗壁；两张湿玻璃片叠在一起就很难分开；胶带能粘纸箱；刚擦完的黑板，板擦上会沾一层粉笔灰。

第二类是“聚”：露珠成球；水黾站在水面上不沉，只踩出四个小凹坑；小心放一枚回形针，它居然能浮在水面；水龙头的水流越往下越细，好像被什么东西往中间收拢。

第三类是“黏”：蜂蜜、糖浆、机油倒起来慢吞吞的，水和酒精却一泻而下。温度一低，蜂蜜黏得更厉害，几乎不肯流动。

第四类是“逃”：湿衣服晾一晾就干；酒精擦在手背上，凉飕飕的，一会儿就没了；香水自己不声不响地“飞”满整个房间。

\begin{shiyan}
\textbf{硬币能盛多少滴水？}

材料：一枚一元硬币、一根滴管（或用筷子尖代替）、清水、一滴洗洁精。

步骤：把硬币平放在桌上，用滴管往硬币正面滴水。一滴、两滴、三滴……一边滴一边数，观察水面鼓成什么样子。滴到水终于漫出来为止，记下滴数。然后在剩下的水里加一小滴洗洁精，搅匀，重新擦干硬币再数一次。

观察点：第一次，水在硬币上堆成一座小小的“水包”，明显高出硬币边缘却不溢出；加了洗洁精之后，同样的硬币只能盛原来三分之一甚至更少的滴数，水面塌塌的。

安全提示：只用常温水和一滴洗洁精，没有危险，做完记得把桌面擦干。
\end{shiyan}

洗洁精做了什么？它既没有把水变稀，也没有把硬币变小。它破坏的是水分子之间的某种“拉拢”。这种拉拢是这一章的主角。

\section{每个分子都被邻居拽着}

把镜头推进到水面。我们已经知道，水和酒精都是由分子组成的，而这些分子从不停歇地做着无规则运动，互相碰撞、乱挤乱撞。温度，从粒子的角度看，就是分子平均运动快慢的量度：温度越高，分子平均跑得越猛。

现在想象一个待在液体内部的水分子。它四面八方都是邻居，前后左右的碰撞大体抵消，它就在原地附近挤来挤去。再看一个恰好停在\textbf{表面}的水分子：它的上面是空气，几乎空无一人；下面和旁边全是同伴。如果这个分子碰巧获得了足够大的向上速度，它就会冲出液面，从此成为水蒸气的一员——这就是蒸发。空气里乱飞的水蒸气分子也可能一头撞回液面，被重新“扣留”——这就是凝结。湿衣服晾干，就是逃走的比回来的多。

\begin{center}
\begin{tikzpicture}[scale=1]
\draw[thick] (0,0) -- (7.6,0);
\foreach \x/\y in {0.6/-0.5, 1.7/-0.9, 2.7/-0.4, 3.7/-0.85, 4.8/-0.45, 5.9/-0.9, 6.9/-0.4}
  {\fill (\x,\y) circle (0.13);}
\foreach \x in {1.1, 2.2, 3.3, 4.4, 5.5, 6.6}
  {\fill (\x,0) circle (0.13);}
\draw[->,thick] (3.3,0.2) -- (3.3,1.3);
\fill (3.3,1.5) circle (0.13);
\draw[->,thick] (5.5,1.4) -- (5.5,0.3);
\fill (5.5,1.6) circle (0.13);
\node at (7.1,-0.7) {\small 液体};
\node at (3.3,1.85) {\small 逃走};
\node at (5.5,1.95) {\small 回来};
\end{tikzpicture}
\end{center}

这幅图是人为的表示法：圆点代表分子，大小、颜色都不代表分子的真实样子，箭头只表示运动方向。

看懂蒸发，就懂了酒精擦在手背上为什么凉：跑得最快、能量最高的那批分子优先逃走，留下的分子平均速度下降——温度降低，你的皮肤就要向它输送热量。同样的道理，狗狗伸舌头、你出汗，都是身体在利用“蒸发带走能量”给自己降温（第 27 章还会从能量的角度再算这笔账）。

蒸发说明：液体里的分子彼此拽着。如果没有拖拽，所有液体都会像泼出去的气体一样瞬间散尽。可这个“拽”是什么？它不是化学键——第 17 章说过，成键断键对应分子本身的改造，而蒸发出来的水分子还是水分子，遇冷又能变回水，全程没有任何一个水分子被拆开。分子之间存在着一种\textbf{比化学键弱得多、却真实存在的吸引}。它弱到一撕就开（想想保鲜膜），又强到足以让亿亿个分子抱成一团，成为液体和固体。

\section{三种“牵手”方式}

分子间的吸引不是一种，而是三种强度不同的“牵手”方式。它们本质上都是电荷之间的电磁吸引，只是电荷分布的花样不同。

\subsection{色散力：连“绝缘户”也逃不掉}

先看一个看似矛盾的事实：氦原子最外层电子住满，谁都不理（第 10 章），甲烷分子 \chem{CH_4} 不带极性，可氦气冷到 $-269\,^{\circ}\mathrm{C}$ 会变成液体，甲烷冷到 $-161\,^{\circ}\mathrm{C}$ 也会液化。\textbf{任何原子、任何分子之间都有吸引，无一例外。}

这种吸引的起因藏在电子身上。电子云不是静止的棉花糖，它在永不停歇地晃动。某一瞬间，电子可能恰好偏向分子的左边，这一瞬间左边就偏负、右边偏正——一个转瞬即逝的“临时电极”。这个临时电极会诱使隔壁分子的电子跟着晃动，两个瞬时偶极互相吸引。下一秒方向变了，吸引依然发生。这种由电子云瞬时涨落引起的吸引叫\textbf{色散力}（也叫伦敦力）。它单看极其微弱，但从不缺席。

色散力有一条非常实用的规律：\textbf{电子越多、分子越大，电子云越容易被晃动，色散力越强}。看证据——卤素单质一家：

\begin{table}[htbp]
\centering
\begin{tabular}{lccc}
\toprule
物质 & 每个分子的电子数 & 沸点 & 室温下（$25\,^{\circ}\mathrm{C}$）的状态 \\
\midrule
氟气 \chem{F_2} & 18 & $-188\,^{\circ}\mathrm{C}$ & 气体 \\
氯气 \chem{Cl_2} & 34 & $-34\,^{\circ}\mathrm{C}$ & 气体 \\
溴 \chem{Br_2} & 70 & $59\,^{\circ}\mathrm{C}$ & 液体 \\
碘 \chem{I_2} & 106 & $184\,^{\circ}\mathrm{C}$ & 固体 \\
\bottomrule
\end{tabular}
\end{table}

四个分子的结构完全一样——都是两个相同原子拉手，都没有永久极性。唯一的差别是电子越来越多。结果从气体一路变成固体：色散力随着电子数的增加稳步增强，强到足以在室温下把碘分子按在固体里。稀有气体也排着同样的队：氦、氖、氩、氪、氙，沸点一路从 $-269\,^{\circ}\mathrm{C}$ 爬到 $-108\,^{\circ}\mathrm{C}$。

再看一个生活例子：天然气里的甲烷（一个碳）是气体，打火机里的丁烷（四个碳）稍一加压就成液体，蜡烛里的石蜡（二三十个碳）在室温下干脆是固体。全是碳氢链，全靠色散力，差别只在分子大小——你的保鲜膜之所以能“贴住”，靠的也是它那些长链分子的色散力，只是被做得非常弱、非常均匀，才能一撕就开。

\subsection{偶极—偶极：带永久“南北极”的分子}

第 20 章讲过，水分子是折线形，氧端偏负、氢端偏正，像一块小磁铁；而二氧化碳虽然每根键都有极性，对称的直线形状却让两极抵消，整体不显极性。像水这样自带永久“正负两极”的分子，会彼此排队：你的正端对我的负端，我的正端对下一个的负端。这种永久偶极之间的吸引叫\textbf{偶极—偶极作用}，一般比色散力强。

一个直接对比：丙酮（洗甲水的主要成分）和丁烷的分子大小、电子数几乎一样，色散力相当；但丙酮分子带永久偶极，沸点是 $56\,^{\circ}\mathrm{C}$，丁烷只有 $-0.5\,^{\circ}\mathrm{C}$。多出来的那几十度，就是偶极—偶极作用“加码”的结果。

\subsection{氢键：偶极作用里的“特种兵”}

当氢原子连在氮、氧、氟这三种“电负性极强”的原子上（电负性见第 11 章），会发生一件特殊的事：氢唯一的电子几乎被邻居整个拉走，氢核——一个孤零零的质子——半裸着暴露在分子表面，带着很强的正电。它会被另一个分子的氮、氧、氟上富集的电子强烈吸引。这种特别强的定向吸引叫\textbf{氢键}。

名字里带个“键”字，请千万记住：\textbf{氢键不是化学键}，它不断裂、不形成分子，只是分子之间一种较强的拉拢。一个水分子的氧端，可以和相邻水分子的氢端形成氢键——这正是水的一切“反常”的起点（第 23 章会专门展开）。

氢键的威力，看一组最漂亮的对比就够了。水分子 \chem{H_2O} 和甲烷分子 \chem{CH_4} 恰好各有 10 个电子，大小也相近——论色散力，两者几乎一模一样。可甲烷的沸点是 $-161\,^{\circ}\mathrm{C}$，水是 $100\,^{\circ}\mathrm{C}$，整整差了 $261$ 度。同样大小的分子，一个要到南极的南极才能留住液态，一个在你手边就是液体。差别全部来自氢键（外带偶极—偶极作用）：每个水分子能同时和周围三四个邻居拉氢键，甲烷分子一个氢键都拉不了。这就是“水在室温下是液体、甲烷是气体”的全部答案——不是水分子内部有什么特别，而是水分子\textbf{之间}多了一重拉拢。氨 \chem{NH_3}、氟化氢 \chem{HF} 也因为氢键，沸点比同族的“亲戚”们高出一截。

\subsection{把强度放在同一把尺上}

三种分子间作用力和化学键比，到底差多少？比较“拉开一摩尔”所需的能量：

\begin{table}[htbp]
\centering
\begin{tabular}{ll}
\toprule
作用类型 & 大致强度（每摩尔） \\
\midrule
色散力（小分子之间） & 约 $1\sim5\ \mathrm{kJ}$ \\
偶极—偶极作用 & 约 $5\sim20\ \mathrm{kJ}$ \\
氢键（如水分子之间） & 约 $10\sim40\ \mathrm{kJ}$ \\
共价键 \chem{C-C}（第 19 章） & 约 $350\ \mathrm{kJ}$ \\
共价键 \chem{O-H} & 约 $460\ \mathrm{kJ}$ \\
\bottomrule
\end{tabular}
\end{table}

一根 \chem{O-H} 共价键，抵得上二十多个水分子间的氢键。这就是“分子间作用力”和“分子内化学键”的本质区别：\textbf{化学键决定分子是什么，分子间作用力决定一大批分子聚在一起时表现出什么状态}——是气体、液体还是固体，是稀是水还是黏如蜜，是能成露珠还是摊成一片。沸腾、蒸发、熔化、黏附，动的都只是分子之间的“手拉手”，分子本身纹丝不动。

\begin{qiaoliang}
\textbf{蒸干一杯水，要付多少“分手费”？}

来做一个真实的数量级估算。把分子彼此拉开需要能量，这笔能量叫作汽化热。对水来说，在 $25\,^{\circ}\mathrm{C}$ 时每蒸发 \ud{1}{mol} 水（也就是 \ud{18}{g}，一小口）大约需要 \ud{44}{kJ} 的能量——这笔钱主要用来拆散水分子之间的氢键和偶极吸引。

一杯水约 \ud{250}{g}，也就是大约 \ud{14}{mol}。把它全部蒸干需要：
\[
14 \times 44 \approx 616\ \mathrm{kJ}
\]
对比一下：把同一杯水从 $20\,^{\circ}\mathrm{C}$ 烧到 $100\,^{\circ}\mathrm{C}$，只需要 $250 \times 4.2 \times 80 \approx 84\ \mathrm{kJ}$。也就是说，\textbf{让水“全部散伙”花的钱，是把它烧开的七倍多}。这就是为什么火锅汤能咕嘟很久才烧干，也是为什么出一场大汗能带走大量体热——每蒸发一克汗，身体就少付一笔可观的热量。

再看反方向的账：酒精的汽化热约 \ud{39}{kJ/mol}，折算到每克约 \ud{0.85}{kJ}，不到水的一半（水约 \ud{2.4}{kJ/g}）——酒精分子之间的拉拢更弱，所以酒精挥发得更快、擦在身上更凉。
\end{qiaoliang}

\section{给看不见的拉拢画个记号}

分子间作用力看不见摸不着，化学家给它们发明了一套简洁的记号。永久偶极用希腊字母 $\delta$（读作“德尔塔”）加正负号表示“部分电荷”：水分子写成氧端 $\delta^{-}$、氢端 $\delta^{+}$——注意不是完整的离子电荷 $+$、$-$，只是“偏负”“偏正”，这是和第 18 章的离子最关键的区别。氢键则用虚线画在分子之间，例如：
\[
\chem{O-H \cdots O-H}
\]
实线是 \chem{O-H} 共价键，分子内部的事；虚线是氢键，分子之间的事。同一张纸上的两种线，强度差着二十倍，千万别看串了。

这套记号依然是人为的表示法。真实的分子间吸引没有“线”，只是电荷之间随距离快速衰减的静电力；而且每一对分子之间的“虚线”都在以万亿分之一秒的时间尺度断开、重建。虚线画的是统计上反复出现的那种姿态，不是一张静止的照片。

\begin{zhengju}
\textbf{分子间的吸引，最初是从理想气体的“裂缝”里看出来的。}

19 世纪，物理学家手里有一条漂亮的定律：理想气体定律。它的前提假设简单粗暴——气体分子是一个个不占体积、互不理睬的小点。这条定律在常温常压下好使得惊人，可一旦加压、降温，真实气体就开始“不听话”：压强总比预言的小一点，而且压到一定程度，气体居然抱团液化了。如果分子之间真的毫无吸引，压缩只会让它们挤得更凶，凭什么突然凝聚成液体？

1869 年，爱尔兰物理学家安德鲁斯系统测量二氧化碳，发现只要温度高于某个特定值（对 \chem{CO_2} 是 $31\,^{\circ}\mathrm{C}$），无论加多大压强都无法把它液化——存在一个“临界点”。气体和液体的界限，原来不是绝对的。

1873 年，荷兰人范德瓦尔斯在博士论文里把这件事变成了理论。他给理想气体方程只加了两笔修正：分子自己占一点体积（记作 $b$），分子之间存在吸引（记作 $a$）。就这两笔，方程不仅解释了真实气体的偏差，还预言了安德鲁斯刚发现的临界点，甚至能算出每种气体被液化的条件。后来的制冷工业——把空气液化、分离出液氮液氧——走的正是这条路。1910 年，范德瓦尔斯因此获得诺贝尔物理学奖。

但故事没完。范德瓦尔斯的 $a$ 是个“哪来的说不清”的参数，同行完全可以质疑：你只是为了凑数据加的修正项吧？真正的回答等了半个世纪的量子力学：1912 年科瑟姆解释了永久偶极之间的吸引，德拜补充了诱导偶极，1930 年伦敦用量子力学证明了电子云瞬时涨落必然产生吸引——哪怕是完全非极性的原子之间。今天我们把分子间的吸引统称为“范德华力”，再把氢键单独记一笔。这段历史值得记住的地方是：先有一个“为了拟合数据”的修正，后有别人设计实验、发展理论去检验它究竟是不是客观存在——修正项最终被证明是大自然的真面目。
\end{zhengju}

\begin{wuqu}
“水烧开，就是水分子被拆成了氢气和氧气。”——错得很典型。壶嘴冒出的白汽（其实是水蒸气凝结成的小液滴）仍然是 \chem{H_2O}，一遇冷就变回水，成分从头到尾没变。想把水分子真正拆开，需要电解或者几千摄氏度的高温：断开一根 \chem{O-H} 共价键要 \ud{460}{kJ/mol}，而沸腾只是拉开分子间的氢键，每个只要约 \ud{20}{kJ/mol}，差着二十多倍。沸腾时剧烈的气泡里装的是水蒸气，不是氢气氧气（顺便说，水刚受热时先冒出的小气泡倒确实主要是溶解在水里的空气）。一句话：\textbf{沸腾拆的是“手拉手”，不是分子本身}。
\end{wuqu}

\section{这套模型什么时候会失灵}

“三种力 + 强度表”这套账房式的分类，解释物态变化、黏度、表面张力非常顺手，但它有几条明确的边界。

第一，它是近似。三种力本质上都是电荷的电磁作用在不同电荷分布下的统计表现，真实世界里它们的界限是模糊的——把一个相互作用归入哪一类，有时取决于你看问题的精度。

第二，它管的是“分子之间”。离子晶体里是带完整电荷的离子在静电吸引（第 18 章），金属里是离域电子的海洋（第 21 章），那两套账不在本章模型的管辖范围。解释食盐为什么熔点高、铜为什么导电，用范德华力会错得离谱。

第三，它解释不了化学反应。糖为什么能在氧气里燃烧、铁为什么生锈，涉及断键和成键，是第四篇的事。分子间作用力再强，也强不过共价键，它决定“物质以什么状态聚集”，不决定“物质变成什么”。

第四，单个分子对的账算不出“集体行为”。冰为什么比水轻、水为什么在 $4\,^{\circ}\mathrm{C}$ 时密度最大、水的表面张力为什么高得出奇——这些反常不是“一根氢键有多强”能回答的，要看亿万个氢键搭成的整张网络。那正是下一章的内容。

\section{回到保鲜膜、露珠、壁虎和香水}

现在兑现开场的承诺：那四件事，背后是\textbf{同一类原因}——分子间的吸引，只是强度和对象不同。

保鲜膜：它是聚乙烯长链分子做成的薄膜，贴附靠的是色散力（厂家还特意调整了配方让“自粘”恰到好处）。色散力本来就弱，薄膜又软，能大面积贴合，所以轻轻一蒙就贴住；也正因为每一处的吸引都微弱，你一撕，成千上万的微弱拉手同时断开，毫不费力。\textbf{粘得牢和撕得开，是同一种力的两面。}

露珠：水分子被同伴往液体内部拉，表面的水分子“最吃亏”——它们上面没有同伴，被横向和向下的邻居拽着，于是水的表面整体倾向于收缩到最小，这就是表面张力。一小滴水，表面张力完胜重力，于是收成一颗近似的球。你在硬币上滴出的那座“水包”，也是表面张力兜出来的；洗洁精分子一头亲水一头亲油，插进水分子之间拆台，表面张力一垮，水包立刻塌方。

壁虎：它的脚底没有吸盘也没有胶，只有几十万根比头发还细的刚毛，每根刚毛末端又分成几百根铲状的绒毛，小到纳米量级。这些绒毛能大面积地贴进墙面分子旁边，让\textbf{色散力}——最普通、人人都有、连氦原子都逃不掉的那种力——以百万计的规模同时发生。2000 年，美国生物学家奥特姆的团队用微型传感器测量了单根刚毛的黏附力，再乘以刚毛数量：理论上足以吊起比壁虎体重重得多的东西，壁虎实际只动用一小部分。它抬脚时把刚毛像撕保鲜膜一样逐根“剥离”，所以走得飞快——你看，连原理都和保鲜膜一样。

香水：香料分子和酒精分子之间的拉拢比较弱，跑得快的分子不断逃进空气，飘来飘去撞进你的鼻腔。瓶子里的液体变少，就是这场旷日持久的“分子大逃亡”的账单。

四种现象，一张底牌。这就是物理化学最迷人的地方：解释力不在于记住四个故事，而在于认出它们是同一个故事。

\section{从壁虎脚到仿生胶带}

理解了分子间作用力，人类开始反过来“雇佣”它。仿照壁虎脚掌，工程师做出了表面布满微米级绒毛的仿生胶带：不用胶水，能反复粘贴几百次，还能在光滑的玻璃上承重。医院里一些贴在皮肤上的传感器贴片、可反复撕贴的医用敷料，思路也来自这里——靠大面积微弱的范德华力，而不是刺激性的化学胶。

黏度的故事也有产业价值：润滑油的核心指标就是黏度，工程师挑选不同长度的碳氢链分子，靠的正是“链越长、色散力越强、越黏稠”这条规律，让发动机在寒冬能启动、在酷暑不流失。

还有一个你绝对想不到的用途：吸入式麻醉药。麻醉气体分子要溶进神经细胞的脂质膜里起效，靠的同样是色散力和偶极作用的强弱搭配——药性差异的底层，是分子间“牵手”的松紧。

不过，这一章我们还欠着一个最大的悬念没还。前面说过，水分子之间靠氢键拉拢——可如果你只知道“拉拢”，水应该是这样的：沸点比甲烷还低，结冰后沉到水底，比热平平无奇。真实的水偏偏条条相反：沸点高得离谱，冰浮在水上，一杯水能“吞”下巨量热量。正是这些反常，让地球有了海洋、气候和体温稳定的身体。为什么？因为氢键不是一根一根孤立地拉，而是亿万分子搭成了一张不断重构的\textbf{网络}。第 23 章，我们专门讲水——讲这张网络怎样让一种最普通的分子，变成生命最不可能缺席的舞台。

\begin{tiaozhan}
范德华力随距离衰减得有多快？物理学家常用一个简化模型——兰纳-琼斯势——来描述两个中性粒子之间的势能：
\[
U(r) = 4\epsilon\left[\left(\frac{\sigma}{r}\right)^{12} - \left(\frac{\sigma}{r}\right)^{6}\right]
\]
其中 $r$ 是两个粒子的间距，$\epsilon$ 和 $\sigma$ 是随物质而定的参数。减号那一项按 $r^{-6}$ 衰减，代表色散吸引：距离加倍，吸引只剩约 $1/64$——这就是为什么保鲜膜必须贴得足够近才“有感觉”。加号那一项按 $r^{-12}$ 变化，代表电子云重叠时急剧上升的排斥：分子靠得太近会被狠狠推开。吸引和排斥在某个距离恰好平衡，液体和固体的“分子间距”就停在这个势能谷底。伦敦在 1930 年证明，$r^{-6}$ 这个幂次不是凑出来的，而是量子力学里电子涨落关联的必然结果。本框涉及一点指数运算和势能曲线，跳过不影响主线；想继续，可以回去对照第 17 章的势能曲线，比较“分子内成键”和“分子间吸引”的谷底一深一浅。
\end{tiaozhan}

\section*{练一练}

\begin{sikaoti}
\item 画一幅粒子示意图：液面上一个分子正在逃逸、另一个分子正在撞回液面，液面下的分子彼此用小弧线表示“拉着”。在图旁标注：哪一层表示蒸发、哪一层表示凝结，并注明这是人为表示法。
\item 乙烷 \chem{C_2H_6}（18 个电子）和甲烷 \chem{CH_4}（10 个电子）都是非极性分子。只用本章的规律预测：谁的沸点更高？说出你依据的那条规律。
\item 蜂蜜比水“黏”得多。用粒子的语言解释：什么叫液体“黏”？（提示：流动时分子要互相滑过去。）再预测：把蜂蜜放进冰箱，黏度会怎样变化，为什么？
\item 查本章的强度表，比较“拉开一个水分子间的氢键”和“断开一根 \chem{O-H} 共价键”所需的能量。用这个对比向一位同学解释：为什么烧水烧一百年也得不到氢气。
\item 画一张能量柱状图：横轴依次是色散力、偶极—偶极作用、氢键、\chem{C-C} 共价键，纵轴用对数或截断的方式表示数量级差异。从图上直接读出“化学键比分子间作用力大约强多少倍”。
\item 露珠在荷叶上缩成滚圆的球，在玻璃窗上却摊成一片。用“水分子之间的拉拢”和“水与固体表面的拉拢”两股力，猜一猜：荷叶表面对水的吸引是强还是弱？这个现象和第 23 章要讲的“相似相溶”有什么关系？
\end{sikaoti}
